Paper~10~\cite{paper10adaptive} introduced a one-threshold magnitude
change-point detector for adaptive recalibration of HygienePrio,
evaluated at the single pre-registered $\tau = 0.05$ with a mixed
outcome (H1 supported, H2 rejected). Paper~10's discussion argued the
H1 success and H2 cost are not separable and are parameterised by
$\tau$. This paper tests that claim directly via a pre-registered
sweep over $\tau \in \{0.02, 0.03, 0.05, 0.075, 0.10\}$ on the
same $K \in \{50, 100, 200\}$, $\lambda = 3$ grid, producing 11{,}250
result rows.

\textbf{The central finding (H3 rejected) is the headline of this
paper}: \textit{no $\tau$ in the swept grid simultaneously satisfies
both Paper~10 pre-registered tolerances}. The worst $K = 200$ adaptive--fixed
delta is exactly $0.0$ at every $\tau$ (the detector either fires or
the unfired-window lag-1 fit happens to match the fixed weights at
this capacity), so the H1 hazard is uniformly avoided. But the
$K = 50$ adaptive--lag1 cell-mean delta is between $-0.041$ and
$-0.053$ at every $\tau$ tested, always exceeding the Paper~10
$-0.02$ pre-registered tolerance. The Pareto curve over $\tau$
sits entirely outside the joint feasible region defined by Paper~10's
two tolerances. The simple magnitude detector is fundamentally
insufficient and richer detectors (CUSUM, Bayesian) are necessary.

\textbf{Two secondary findings.} (1) H1's worst-case hazard is
uniformly zero across the entire $\tau$ grid --- a strictly stronger
result than Paper~10's single-$\tau$ H1. The simple detector
\textit{robustly} avoids the K=200 hazard. (2) H2 and H4 were
pre-registered with the wrong sign direction: the protocol predicted
fire fraction would be monotone non-decreasing in $\tau$ (more firing
at higher threshold) but the actual relationship is strongly
negative ($\rho = -0.97$, $p = 0.005$). Larger $\tau$ requires a
larger shift to trigger and therefore fires less. The H2 prediction
inherited this sign error and is similarly rejected with reverse sign
($\rho = +0.87$). Both rejections are reported honestly per the
protocol's stop rules; the data direction is the mathematically
correct one, and we use it as the discussion's framing.

The deployable conclusion is sharp: \textbf{the simple magnitude
detector is uniformly safe at $K = 200$ across the studied $\tau$
range but uniformly pays an unavoidable $\geq 4$~pp penalty at
$K = 50$.} The operating-point question Paper~10 left open has a
negative answer for this detector family; the answer requires a
richer detector that has not been built. All results are bounded to
the synthetic EEHDA context.
